\documentclass[12pt]{article}
\usepackage[utf8]{inputenc}
\usepackage[T1]{fontenc}
\usepackage[margin=1in]{geometry}
\usepackage{setspace}
\usepackage{xcolor}
\usepackage[colorlinks=true, urlcolor=blue!60!black, linkcolor=blue!60!black]{hyperref}
\setlength{\parindent}{0em}
\setlength{\parskip}{0.6em}
\renewcommand{\baselinestretch}{1.0}
\begin{document}
\begin{center}
{\LARGE \textit{Hammer v. Dagenhart}, 247 U.S. 251 (1918)\par}
\vspace{0.5em}
{\large Case Brief --- Two-Page Edition\par}
\end{center}

\section{Memory Jogger}
Congress could not bar interstate shipment of harmless goods made with child labor; the dual-federalism holding was overruled in \textit{United States v. Darby} (1941) (\href{https://supreme.justia.com/cases/federal/us/247/251/}{247 U.S. 251}).

\section{Detailed Case Facts}
The Keating-Owen Child Labor Act of 1916 barred interstate shipment of goods from mines employing children under sixteen and from factories employing children under fourteen, or children of fourteen to sixteen for long hours or at night. Roland Dagenhart, whose two minor sons worked with him in a Charlotte cotton mill, sued U.S. Attorney W. C. Hammer to enjoin enforcement.

\section{Procedural History}
The district court held the Act unconstitutional and enjoined it; Hammer appealed directly to the Supreme Court, which affirmed.

\section{Judicial Votes}
Affirmed, 5--4: Day wrote for White, Van Devanter, Pitney, and McReynolds; Holmes dissented, joined by McKenna, Brandeis, and Clarke.

\section{Holding}
Congress may not use the commerce power to exclude intrinsically harmless goods because of the local labor conditions of their production. Manufacturing is not commerce, and regulating it invades the police power reserved to the States by the Tenth Amendment.

\section{Analysis of Opinions}
Day distinguished the lottery, impure-food, prostitution, and liquor cases because there ``the use of interstate transportation was necessary to the accomplishment of harmful results.'' The cloth was harmless and the harm occurred before shipment; Congress also had no power to equalize competition among states. Holmes replied that the Act regulated interstate shipment itself, that ``regulation means the prohibition of something,'' and that a valid federal power is not void because of its indirect effects on state policy.

\section{Examples: Future Applications}
\textbf{Same-Side:} a ban on shipping furniture made below a federal minimum wage would fall, because the goods are harmless and the labor local. \textbf{Opposite-Side:} a ban on shipping toxic toys or stolen cars stands, because transportation spreads or completes the harm. \textbf{Fence-Sitter:} a ban on prison-made goods only into states that forbid them, later upheld in \textit{Kentucky Whip \& Collar}.

\section{Critique}
Textualists fault Day's insertion of ``expressly'' into the Tenth Amendment, though some originalists defend the commerce--manufacture line. Progressives call the decision formalist and blind to the race to the bottom that made national regulation necessary, the rationale of \textit{Darby}.

\section{Key Quotations}
\begin{quote}
``The making of goods and the mining of coal are not commerce, nor does the fact that these things are to be afterwards shipped or used in interstate commerce make their production a part thereof.''
\end{quote}
\noindent --- Justice Day, Opinion of the Court, 247 U.S. at 272.

\end{document}
